\documentclass[conference]{IEEEtran}

\usepackage{amsmath}
\usepackage{amssymb}
\usepackage{graphicx}
\usepackage{cite}

\title{Periodic Optimal Dynamic Pricing against Online-Learning Users in Smart Grid Demand Response}

\author{\IEEEauthorblockN{Author Name}
\IEEEauthorblockA{School of \ldots\\
University\\City, Country\\email@example.com}}

\begin{document}

\maketitle

\begin{abstract}
This paper studies how an electricity aggregator should set dynamic prices when users update their demand strategies with the Hedge (multiplicative weights) learning algorithm. We model the repeated interaction as a finite-horizon dynamic game, formulate the aggregator's problem as an optimal control problem over a state transition graph, and show that the optimal price sequence eventually becomes periodic under commensurable payoff parameters. We compare periodic pricing with static Stackelberg pricing and myopic pricing through numerical experiments.
\end{abstract}

\begin{IEEEkeywords}
demand response, dynamic pricing, Hedge algorithm, repeated games, Stackelberg game
\end{IEEEkeywords}

\section{Introduction}
\IEEEPARstart{T}{he} increasing deployment of smart meters and automated demand management raises a fundamental question for electricity pricing: how should an aggregator price electricity when users are not fully rational, but instead adapt their behavior through online learning algorithms?

The optimal pricing formula for rational users was recently derived in a Stackelberg game framework \cite{mu2024pricing}. In parallel, the optimal strategy against the Hedge algorithm was characterized for repeated zero-sum games \cite{guo2023optimal}. This paper connects these two lines: we study dynamic pricing when users are Hedge learners.

Our main contributions are:
\begin{itemize}
    \item a repeated-game model of dynamic pricing against Hedge-learning users;
    \item a state-transition formulation and Bellman equation for the aggregator;
    \item structural evidence that optimal pricing becomes periodic;
    \item a numerical comparison with static Stackelberg and myopic baselines.
\end{itemize}

\section{Problem Formulation}
Consider a finite horizon $T$. At each stage $t \in \{1,\dots,T\}$, the aggregator chooses a price $p_t$ from a finite set $\mathcal{P}$, and each user $i \in \{1,\dots,N\}$ chooses a demand level from a finite set $\mathcal{D} = \{d_1, \dots, d_K\}$.

\subsection{User Utility}
User $i$ receives utility
\begin{equation}
    u_i(d, p) = v_i(d) - p d,
\end{equation}
where $v_i(d)$ is the value of consuming level $d$. We assume that $v_i$ is increasing and concave in $d$.

\subsection{Hedge Users}
Users do not solve the stage game myopically. Instead, each user updates a multiplicative weight for every demand action. The weights evolve as
\begin{equation}
    w_{i,a,t+1} = w_{i,a,t} \exp\bigl(\eta u_i(a, p_t)\bigr),
\end{equation}
where $\eta$ is the learning rate. The stage strategy is
\begin{equation}
    x_{i,a,t} = \frac{w_{i,a,t}}{\sum_{b \in \mathcal{D}} w_{i,b,t}}.
\end{equation}
The initial weights are $w_{i,a,1} = 1$.

\subsection{Aggregator Objective}
The aggregator's stage revenue is
\begin{equation}
    R_t = p_t \sum_{i=1}^{N} \sum_{a \in \mathcal{D}} x_{i,a,t} d_a,
\end{equation}
and its objective is to maximize the cumulative expected revenue
\begin{equation}
    \max_{p_1,\dots,p_T \in \mathcal{P}} \sum_{t=1}^{T} R_t.
\end{equation}

The aggregator knows $\eta$ and the users' utility functions, but users are not treated as rational Stackelberg followers.

\section{Dynamics and Bellman Equation}
For a user with two demand levels, the weight ratio
\begin{equation}
    \frac{w_{i,1,t}}{w_{i,2,t}} = \exp(\eta s_{i,t})
\end{equation}
defines a scalar state $s_{i,t}$. The state evolves as
\begin{equation}
    s_{i,t+1} = s_{i,t} + \Delta_i(p_t),
\qquad
    \Delta_i(p) = u_i(d_1, p) - u_i(d_2, p).
\end{equation}
When users are homogeneous, the aggregate state is a single scalar $s_t$, and the expected demand is
\begin{equation}
    D_t(p, s) = N \left( \frac{d_1 e^{\eta s_t} + d_2}{e^{\eta s_t} + 1} \right).
\end{equation}

The aggregator's problem is then a finite-horizon optimal control problem:
\begin{equation}
    V_t(s) = \max_{p \in \mathcal{P}} \Bigl[ p D_t(p, s) + V_{t+1}\bigl(s + \Delta(p)\bigr) \Bigr],
\qquad
    V_{T+1}(s) = 0.
\end{equation}

Direct dynamic programming over the state grid has complexity $O(T |\mathcal{S}| |\mathcal{P}|)$, which grows with the horizon. This motivates the study of periodicity in the optimal price sequence.

\section{Structural Properties and Algorithm}
We aim to establish the following results:
\begin{enumerate}
    \item \textbf{Static pricing is not optimal.} A constant Stackelberg price exploits rational users, but against Hedge learners the aggregator can improve revenue by steering the user's learning dynamics.
    \item \textbf{Eventual periodicity.} When all state steps $\Delta(p)$ are commensurable, i.e., $\Delta(p) \in \mathbb{Q}$ for every $p \in \mathcal{P}$, the optimal price sequence eventually enters a cycle.
    \item \textbf{Efficient computation.} Once the transient and the period are identified, dynamic programming needs to be solved only over one period plus a short terminal segment.
\end{enumerate}

The planned algorithm is:
\begin{enumerate}
    \item compute the myopic path and the zero path;
    \item detect the transient length and the period length;
    \item solve the Bellman equation over one period;
    \item repeat the periodic segment and solve the tail directly.
\end{enumerate}

\section{Numerical Experiments}
We plan to compare four pricing rules:
\begin{itemize}
    \item static Stackelberg pricing;
    \item myopic best-response pricing;
    \item a periodic heuristic;
    \item the dynamic programming solution over the periodic state graph.
\end{itemize}
For each rule we report cumulative revenue, average revenue per period, transient length, and period length. The simulation code is available in the companion repository.

\section{Conclusion}
This paper models dynamic electricity pricing against Hedge-learning users as a repeated game and reduces the aggregator's problem to a Bellman equation on a scalar state. The key insight is that optimal pricing can become periodic, which makes the problem computationally tractable and economically exploitable. Future work includes decaying learning rates, unknown learning algorithms, heterogeneous users, and storage-aware demand response.

\bibliographystyle{IEEEtran}
\bibliography{refs}

\end{document}
